\section{A superlogarithmic lower bound}\label{sec:lower}

We now give a lower bound that does not use
Theorem~\ref{thm:main}. The construction uses many primes that can
cancel simultaneously in the symbolic denominator. The key is to
arrange these cancellations at enough different period lengths, not
just at one height.

\begin{proof}[Proof of Theorem~\ref{thm:lower}]
Write \(F_m=2^{-m}P(3^m)\), where \(P(Y)\) is the pure count defined
in Section~\ref{sec:assembly}. By \eqref{eq:normalized-preperiod}, to bound even \(N_C(T)-P(T)\) from below
it suffices to bound the average of the \(F_m\) from below.

\smallskip
\noindent\emph{A common multiple of many prime orders.}
Fix \(\beta>15/(32\sqrt e)\) with \(\beta<\kappa\).
By Lichtman~\cite[Theorem~1.1]{Lichtman2022}, extending
Baker and Harman~\cite[Theorem~1]{BH1998}, for every sufficiently large
integer \(x\) there are \(L\gg_\beta x/(\log x)^{C_\beta}\) primes
in \((x,2x]\) whose predecessors have no prime factor exceeding
\(x^\beta\), for a fixed \(C_\beta>1\). Take \(\mathcal L\) to be the
set of these primes and put \(Y=x\). Write
\[
 Q=\prod_{p\in\mathcal L}p,\qquad
 R=x^4\prod_{\substack{r\le x^\beta\\r\ {\rm prime}}}
                       r^{\lfloor\log(2x)/\log r\rfloor}.
\]
Every \(p-1\) with \(p\in\mathcal L\) divides \(R\): each of its
prime powers is at most \(2x\), and its prime base is at most
\(x^\beta\). For large \(x\) all the selected primes exceed \(3\).
Fermat's little theorem then gives \(Q\mid 3^n-1\) whenever
\(n\) is a positive multiple of \(R\).
We have \(\log Q\le 2x\log(2x)\) and
\(\log R\ll x^\beta\), the latter from
\(\log R\le4\log x+\pi(x^\beta)\log(2x)\).
Here the elementary Chebyshev bound \(\pi(z)\ll z/\log z\) suffices:
\(\prod_{m<r\le2m, r\ {\rm prime}}r\) divides \(\binom{2m}{m}\),
so dyadic summation gives \(\sum_{r\le z, r\ {\rm prime}}\log r\ll z\).
Separating primes at \(\sqrt z\) gives the claimed bound for \(\pi(z)\).

\smallskip
\noindent\emph{Many cancelled words at each selected period.}
Set
\[
 k=\left\lfloor\frac{Y^\kappa L}{8Y}\right\rfloor,
 \qquad j=\lfloor k\log_3Y\rfloor .
\]
Since \(Y=x\), we have
\(k\gg x^\kappa/(\log x)^{C_\beta}\) for large \(x\). Also \(L\le2Y\),
so \(k\ll x^\kappa\), \(1\le k<L\), and \(j\ll x^\kappa\log x<R\).

For each positive multiple \(n\) of \(R\) let
\[
 X_n=\sum_{t=0}^{n-1}b_t3^t,
\]
where the \(b_t\) are independent fair bits. Thus \(2X_n\) is
uniform among the integers specified by the \(2^n\) length-\(n\)
Cantor words. To compare \(X_n\bmod Q\) with the uniform residue,
consider a nonzero frequency \(s\bmod Q\). Its Fourier coefficient
has modulus
\[
 \left|\mathbb E\exp(2\pi i sX_n/Q)\right|
   =\prod_{t=0}^{n-1}|\cos(\pi s3^t/Q)|.
\]
Let \(\vartheta=\cos(\pi/6)<1\) and
\(B_Q=\lceil\log_3 Q\rceil+2\).
Every block of \(B_Q\) consecutive factors contains one factor at
most \(\vartheta\). Indeed the distance from \(s3^t/Q\) to the
nearest integer is at least \(1/Q\), since \(3\) is invertible
modulo \(Q\) and \(s\ne0\). Whenever that distance is less than
\(1/6\), tripling the phase triples the distance. It cannot stay
below \(1/6\) for the whole block. The Fourier coefficient is
therefore at most \(\vartheta^{\lfloor n/B_Q\rfloor}\).
By character orthogonality, at each single residue the error from
uniform probability is at most this bound. For any residue event
the error is thus at most
\[
 Q\vartheta^{\lfloor n/B_Q\rfloor}
 \le \exp\bigl(-c_0x^3/\log x\bigr)
\]
for an absolute \(c_0>0\) and large \(x\), because \(n\ge R\ge x^4\)
and \(\log Q\le 2x\log(2x)\). This bound is uniform in \(n\).

For a uniform residue modulo \(Q\), the events \(p\mid X_n\),
\(p\in\mathcal L\), are independent by the Chinese remainder
theorem. The probability of exactly \(k\) successes in that model
is at least
\[
 \binom Lk(2Y)^{-k}(1-1/Y)^L
   \gg \left(\frac{L}{2Yk}\right)^k,
\]
since \(\binom Lk\ge(L/k)^k\) and \(L\le2Y\).
This lower bound is at least \(\exp(-O(x^\kappa\log x))\),
using \(L/(2Yk)\ge1/(2Y)\). The uniform Fourier error is
negligible compared with it. It follows that at least
a constant times \(2^n(L/(2Yk))^k\) Cantor words have exactly
\(k\) selected primes dividing \(X_n\).

Each such word represents the point \(2X_n/(3^n-1)\).
The common divisor of its displayed numerator and denominator
is at least \(Y^k\ge3^j\); its actual reduced denominator is
therefore at most \(3^{n-j}\) and is prime to \(3\).
For fixed \(n\) the fractions are distinct because the numerators
at the common denominator \(3^n-1\) are distinct. Shorter least
periods cause no duplication at this fixed \(n\), and we do not
count points from different lengths as disjoint. Thus
\[
 F_{n-j}\gg 2^j\left(\frac{L}{2Yk}\right)^k
      \ge \frac12\left(\frac{Y^\kappa L}{2Yk}\right)^k
      \ge \tfrac12\,4^k .
\]
The indices \(n-j\) form a progression of positive integers with
spacing \(R\), starting at \(R-j\).

\smallskip
\noindent\emph{Averaging over the progression.}
For \(M\ge4R\) there are at least \(M/(2R)\) such indices less
than \(M\). Consequently there is an absolute \(c_1>0\) such that
\[
 \frac1M\sum_{m<M}F_m\gg \frac{4^k}{R}
     \ge\exp\left(c_1\frac{x^\kappa}{(\log x)^{C_\beta}}\right)
\]
for large \(x\), because \(\log R=O(x^\beta)\) and \(\beta<\kappa\).
For each sufficiently large \(M\), take
\(x=\lfloor a(\log M)^{1/\beta}\rfloor\) for a sufficiently small
fixed \(a>0\). The bound \(\log R\ll x^\beta\) ensures \(R\le M^{1/2}\),
so \(M\ge4R\). The lower bound for the average becomes
\[
 \exp\left(c_2\frac{(\log M)^{\kappa/\beta}}
                        {(\log\log M)^{C_\beta}}\right)
\]
up to a constant depending on \(\beta\), with \(c_2>0\). Inserting this in
the preperiod identity bounds \((N_C(3^M)-P(3^M))/(2^M M)\) from below.

For arbitrary large \(T\), put \(M=\lfloor\log_3T\rfloor\).
The numerator \(N_C(T)-P(T)\) counts the points of positive denominator
valuation at \(3\), so it is nondecreasing. Also
\(T^\kappa<2\cdot 2^M\) and \(M\asymp\log T\). Hence the argument gives
constants \(c_\beta>0\) and \(C_\beta>1\) such that
\[
 N_C(T)-P(T)\gg_\beta T^\kappa\log T
 \exp\left(c_\beta\frac{(\log\log T)^{\kappa/\beta}}
                             {(\log\log\log T)^{C_\beta}}\right)
\]
for all sufficiently large \(T\). For every \(0<\rho<\rho_*\), choose
\(\beta>15/(32\sqrt e)\) with \(\beta<\kappa\) and \(\rho<\kappa/\beta\).
The displayed exponent then eventually exceeds \((\log\log T)^\rho\)
and proves the theorem.
\end{proof}

The argument also shows that the averages
\(M^{-1}\sum_{m<M}F_m\) tend to infinity, rather than being large
only along a subsequence. The lower bound does not determine the
optimal subpower factor in \(N_C(T)\).

\begin{proof}[Proof of Theorem~\ref{thm:lower-cylinders}]
Write \(A\) for the integer of the specified prefix. The cylinder map
is \(y\mapsto (A+y)/3^d\), a bijection from \(C\) onto \(I\).
If \(y=a/q\) is reduced and \(3\mid q\), then \(qA+a\) is coprime to
\(q\) and to \(3\), so the image has reduced denominator exactly
\(3^d q\). If instead \(3\nmid q\), the image has denominator
valuation at most \(d\). This proves the identity at every height,
including the endpoints. The lower bound for \(N_C(Y)-P(Y)\) proved
above gives the first inequality. At triadic heights \(T=3^M\)
with \(M\ge d+1\), the same identity reads
\[
 2^{M-d-1}\sum_{0\le m<M-d}F_m .
\]
For \(d\le\theta\lfloor\log_3 T\rfloor\), we have
\(Y\ge T^{1-\theta}\), \(Y^\kappa=2^{-d}T^\kappa\) and
\(\log\log Y=\log\log T+O_\theta(1)\). Apply the first inequality
with any fixed exponent strictly between \(\rho\) and \(\rho_*\)
to absorb the last bounded change of scale. This gives the claimed
uniform estimate.
\end{proof}
